\chapter{Relevância do Trabalho}\label{cap:relevancia}

As doenças hepáticas representam uma significativa carga global de saúde, sendo responsáveis por elevado número de internações, complicações graves e mortalidade, a literatura aponta para a necessidade de transição da prática tradicional para a medicina de precisão em hepatologia, graças ao volume crescente de dados clínicos, laboratoriais, de imagem e genômicos disponíveis \citep{SABRIE2023350}. No Brasil, a mortalidade por hepatites crônicas aumentou em algumas regiões (Norte e Nordeste) enquanto outras formas de hepatite apresentaram tendência de queda entre 2001 e 2020, reforçando a urgência de ferramentas de monitoramento de evolução clínica \citep{sousa2023}. 

Em paralelo, os prontuários eletrônicos e notas clínico-assistenciais contêm riqueza de dados em texto livre, que até o momento permanecem subutilizados devido à complexidade do processamento linguístico, revisões recentes em hepatologia e gastroenterologia ressaltam que o PLN já consegue extrair informações úteis desses textos, mas ainda padece de limitações como falta de generalização, baixa explicabilidade, escassez de integração temporal e escopo restrito \citep{omar2025}. 

Nesse cenário, a combinação de PLN com grafos de conhecimento e modelos de recuperação assistida, como a abordagem GraphRAG, surge como uma oportunidade inovadora para capturar relações complexas entre entidades clínicas, modelar trajetórias ao longo do tempo e permitir consultas mais ricas e estruturadas, algo ainda raro no contexto de hepatologia \citep{stammers2025}. Portanto, este trabalho se mostra relevante ao propor uma solução que aborda três lacunas principais: (1) a extração inteligente de dados textuais não estruturados em hepatopatias, (2) o monitoramento longitudinal das evoluções clínicas num ambiente seguro e local, e (3) o uso de grafos de conhecimento para dar sentido e facilitar a exploração dessas informações por clínicos, pesquisadores e gestores. Essa relevância é tanto prática — contribuindo para melhoria no acompanhamento, prognóstico e tratamento de pacientes com doenças hepáticas — quanto metodológica — ao aplicar técnicas de ponta em PLN e KG num domínio ainda pouco explorado.  
